\documentclass[tikz,border=6pt]{standalone}
\usepackage[no-math]{fontspec}
\usepackage[brazil]{babel}
\usepackage{amsmath,amssymb}
\IfFileExists{C:/Windows/Fonts/Lato_Regular.ttf}{%
  \setmainfont{Lato}[
    Path = C:/Windows/Fonts/,
    UprightFont = Lato_Regular.ttf,
    BoldFont = Lato_Bold.ttf,
    ItalicFont = Lato_Italic.ttf,
    BoldItalicFont = Lato_BoldItalic.ttf
  ]
  \setsansfont{Lato}[
    Path = C:/Windows/Fonts/,
    UprightFont = Lato_Regular.ttf,
    BoldFont = Lato_Bold.ttf,
    ItalicFont = Lato_Italic.ttf,
    BoldItalicFont = Lato_BoldItalic.ttf
  ]
}{%
  \RequirePackage[default]{lato}
}
\renewcommand{\familydefault}{\sfdefault}

\usepackage{xcolor}
\usepackage{tikz}
\usepackage{pgfplots}
\pgfplotsset{compat=1.18}
\usetikzlibrary{arrows.meta,positioning,shapes.geometric,fit,backgrounds,calc,decorations.pathreplacing,patterns}

% Paleta institucional UEFS / CHF662
\definecolor{uefsNavy}{HTML}{1B2A4A}
\definecolor{uefsTeal}{HTML}{0E7C7B}
\definecolor{uefsForest}{HTML}{2D6A4F}
\definecolor{uefsGold}{HTML}{C4983B}
\definecolor{uefsGoldDark}{HTML}{8F6B1E}
\definecolor{uefsAlert}{HTML}{B02A37}
\definecolor{uefsOrange}{HTML}{D95F02}
\definecolor{uefsInk}{HTML}{22252A}
\definecolor{uefsMuted}{HTML}{5C6470}
\definecolor{uefsWater}{HTML}{2E86AB}
\definecolor{uefsSand}{HTML}{D9B382}
\definecolor{bgNavy}{HTML}{F2F5F9}
\definecolor{bgTeal}{HTML}{E9F5F4}
\definecolor{bgForest}{HTML}{EEF5F0}
\definecolor{bgGold}{HTML}{FCF7EB}
\definecolor{bgAlert}{HTML}{FDF2F2}
\definecolor{bgGray}{HTML}{F8FAFC}
\definecolor{lineGray}{HTML}{CFD8DC}

\tikzset{
  >=Stealth,
  aaCard/.style={draw=#1, fill=white, rounded corners=4pt, line width=1pt, inner sep=6pt, align=center, font=\small},
  aaBox/.style={draw=#1, rounded corners=4pt, line width=1.1pt, inner sep=6pt, align=center, font=\footnotesize\sffamily, text=uefsInk},
  aaTitle/.style={font=\normalsize\bfseries\sffamily, text=uefsNavy, align=center},
  aaSub/.style={font=\footnotesize\sffamily, text=uefsMuted, align=center},
  aaArrow/.style={->, >=Stealth, line width=1.4pt},
  aaDArrow/.style={<->, >=Stealth, line width=1.4pt}
}


\begin{document}
\begin{tikzpicture}[
  licBox/.style={
    draw=#1,
    fill=#1!7!white,
    rounded corners=5pt,
    line width=1.3pt,
    text width=6.1cm,
    minimum height=1.85cm,
    inner sep=8pt,
    align=center,
    font=\footnotesize\sffamily,
    text=uefsInk
  }
]

% ================= PAINEL (a): FLUXOGRAMA VERTICAL =================
\node[aaTitle, font=\normalsize\bfseries\sffamily, text=uefsNavy] (titleA) at (3.3, 6.95)
  {(a) Licenciamento Ambiental Trifásico\\(Res. CONAMA 237/1997)};

\node[licBox=uefsNavy] (lp) at (3.3, 5.15) {%
  {\small\bfseries\color{uefsNavy}LP (Licença Prévia)}\\[3pt]
  Aprova localização e viabilidade ambiental\\
  (etapa do EIA/RIMA e audiências)%
};

\node[licBox=uefsGoldDark] (li) at (3.3, 2.35) {%
  {\small\bfseries\color{uefsGoldDark}LI (Licença de Instalação)}\\[3pt]
  Autoriza o início das obras conforme\\
  planos de controle e PBA%
};

\node[licBox=uefsForest] (lo) at (3.3, -0.45) {%
  {\small\bfseries\color{uefsForest}LO (Licença de Operação)}\\[3pt]
  Autoriza funcionamento após verificar\\
  cumprimento das condicionantes%
};

% Setas verticais nítidas e bem visíveis entre as caixas
\draw[aaArrow, color=uefsNavy, line width=1.7pt, shorten <=3pt, shorten >=3pt] (lp.south) -- (li.north);
\draw[aaArrow, color=uefsGoldDark, line width=1.7pt, shorten <=3pt, shorten >=3pt] (li.south) -- (lo.north);

% ================= PAINEL (b): BARRAS HORIZONTAIS PGFPLOTS =================
\begin{axis}[
  at={(10.0cm, -1.35cm)},
  anchor=south west,
  width=9.6cm,
  height=8.4cm,
  xbar,
  bar shift=0pt,
  bar width=0.68cm,
  xmin=0,
  xmax=760,
  ymin=0.35,
  ymax=5.65,
  ytick={1,2,3,4,5},
  yticklabels={< 10 m, 10 a 50 m, 50 a 200 m, 200 a 600 m, > 600 m},
  xtick={0,100,200,300,400,500,600},
  xlabel={Largura mínima da APP a partir da calha regular (m)},
  ylabel={Largura do curso d'água},
  xlabel style={font=\footnotesize\sffamily, text=uefsInk},
  ylabel style={font=\footnotesize\sffamily, text=uefsInk},
  title={(b) APPs Marginais: Código Florestal (Lei nº 12.651/2012)},
  title style={font=\normalsize\bfseries\sffamily, text=uefsNavy, yshift=4pt},
  axis lines*=left,
  axis line style={color=lineGray, line width=0.9pt},
  tick style={color=uefsMuted},
  xticklabel style={font=\footnotesize\sffamily, text=uefsInk},
  yticklabel style={font=\footnotesize\sffamily, text=uefsInk},
  clip=false
]

\addplot[xbar, bar shift=0pt, fill=uefsTeal, draw=none] coordinates {(30, 1) (50, 2)};
\addplot[xbar, bar shift=0pt, fill=uefsForest, draw=none] coordinates {(100, 3) (200, 4)};
\addplot[xbar, bar shift=0pt, fill=uefsNavy, draw=none] coordinates {(500, 5)};

% Rótulos das faixas marginais
\node[right=4pt, font=\footnotesize\bfseries\sffamily, text=uefsInk] at (axis cs:30,  1) {30 m (cada margem)};
\node[right=4pt, font=\footnotesize\bfseries\sffamily, text=uefsInk] at (axis cs:50,  2) {50 m (cada margem)};
\node[right=4pt, font=\footnotesize\bfseries\sffamily, text=uefsInk] at (axis cs:100, 3) {100 m (cada margem)};
\node[right=4pt, font=\footnotesize\bfseries\sffamily, text=uefsInk] at (axis cs:200, 4) {200 m (cada margem)};
\node[right=4pt, font=\footnotesize\bfseries\sffamily, text=uefsInk] at (axis cs:500, 5) {500 m (cada margem)};

\end{axis}

\end{tikzpicture}
\end{document}
